\section{Related Work}
\label{sec:related}

We position CPSForge at the intersection of LLM-driven CPS attack generation, autonomous cyber agents, PLC-based testbeds, and adaptive CPS defense. Table~\ref{tab:comparison} summarizes differences.

\textbf{LLM-based CPS attacks.} Ahmed~\cite{ahmed2025attackllm} introduced AttackLLM, using GPT-4 to generate attack patterns for the SWaT testbed, but without a live execution loop, safety shield, or defender. Gowdanakatte et al.~\cite{gowdanakatte2025threat} proposed LLM-based CPS threat emulation at the planning level without PLC-level execution. Ning et al.~\cite{ning2025malf} presented MALF for multi-agent protocol fuzzing at the network layer (Modbus/S7Comm packets), not process-level tags. Garcia et al.~\cite{garcia2017harvey} introduced HARVEY, a physics-model-aware MitM rootkit for PLC I/O; CPSForge replaces hand-crafted physics models with general-purpose LLM reasoning and adds an autonomous defender.

\textbf{Autonomous cyber agents.} PentestGPT~\cite{deng2024pentestgpt}, PentestAgent~\cite{shen2025pentestagent}, and Hackphyr~\cite{rigaki2024hackphyr} demonstrate LLM-driven red teaming on IT infrastructure. Xu et al.~\cite{xu2025forewarned} identify a significant gap in OT/CPS-specific agents. CPSForge fills this gap with attacker and defender agents on physical process tags within a real PLC loop.

\textbf{PLC-based testbeds.} SWaT~\cite{mathur2016swat, goh2017swat}, ICSSIM~\cite{dehlaghi2023icssim}, SPHERE~\cite{garcia2025sphere}, SIMPLE-ICS~\cite{pramadi2026simpleics}, and domain-specific HIL platforms~\cite{chen2024nuclear_hil, nam2025hintsec, mathew2024hil_thermal} provide the realism needed for CPS security evaluation. Factory~I/O has been used for penetration testing~\cite{lazaridis2023factoryio_pentest} and digital-twin analysis~\cite{wang2024factoryio_dt}. None integrate LLM-driven attack or defense agents.

\textbf{Adaptive CPS defense.} Evo-Defender~\cite{hu2025evodefender} co-evolves attacker/defender via iterative fuzzing on simulators. Lanotte et al.~\cite{lanotte2022runtime_ics} formalize runtime enforcement for ICS; Baird et al.~\cite{baird2024scalable_enforcement} extend it to scalable architectures. LLM4PLC~\cite{fakih2024llm4plc} and Agents4PLC~\cite{liu2026agents4plc} demonstrate LLM reasoning about PLC code. Ann et al.~\cite{ann2024ids_llm} investigate LLM-based ICS intrusion detection. CPSForge combines LLM red/blue teaming with a mandatory safety shield and execution on a real PLC.

\begin{table}[t]
\centering
\caption{Comparison of CPSForge with related systems.}
\label{tab:comparison}
\resizebox{\columnwidth}{!}{%
\begin{tabular}{lccccccc}
\hline
\textbf{System} & \rotatebox{70}{\textbf{LLM Attacker}} & \rotatebox{70}{\textbf{LLM Defender}} & \rotatebox{70}{\textbf{Safety Shield}} & \rotatebox{70}{\textbf{Real PLC}} & \rotatebox{70}{\textbf{Plant Sim.}} & \rotatebox{70}{\textbf{Live Loop}} & \rotatebox{70}{\textbf{Fine-Tune}} \\
\hline
AttackLLM~\cite{ahmed2025attackllm}       & \checkmark &            &            &            & \checkmark &            &            \\
Evo-Defender~\cite{hu2025evodefender}     &            &            &            &            & \checkmark & \checkmark &            \\
MALF~\cite{ning2025malf}                  & \checkmark &            &            &            &            &            &            \\
PentestGPT~\cite{deng2024pentestgpt}      & \checkmark &            &            &            &            & \checkmark &            \\
Hackphyr~\cite{rigaki2024hackphyr}        & \checkmark &            &            &            &            & \checkmark & \checkmark \\
HARVEY~\cite{garcia2017harvey}            &            &            &            &            & \checkmark & \checkmark &            \\
\textbf{CPSForge (ours)}                  & \checkmark & \checkmark & \checkmark & \checkmark & \checkmark & \checkmark & \checkmark \\
\hline
\end{tabular}%
}
\end{table}
